% Vietnamese translation for OI Public Library
% Source-PDF: 比赛 CONTEST/2021“MINIEYE杯”中国大学生算法设计超级联赛/2021“MINIEYE杯”中国大学生算法设计超级联赛（2）/2021“MINIEYE杯”中国大学生算法设计超级联赛（2）-题目集.pdf
\documentclass[11pt]{article}
\usepackage{oipl-vn}

\title{2021 MINIEYE Cup, HDU Multi-University Training Contest 2: Đề bài}
\author{}
\date{}

\begin{document}
\maketitle

\origtitle{2021 MINIEYE Cup Multi-University Training Contest 2, Problem Set}

\section{1001. I love cube}

\paragraph{Giới hạn.} Thời gian: 1 giây. Bộ nhớ: 512 MB.

Cho một hình lập phương có cạnh dài $n-1$.  Hãy đếm số tam giác đều có ba đỉnh
là các điểm nguyên trong hoặc trên biên hình lập phương.  Mỗi cạnh của tam giác
phải song song với một trong ba mặt phẳng tọa độ $Oxy$, $Oxz$, $Oyz$.

\paragraph{Input.}
Dòng đầu chứa số nguyên dương $T$ ($T\le 10^5$).  Sau đó là $T$ bộ test, mỗi
bộ gồm một số nguyên $n$ ($1\le n\le 10^{18}$).

\paragraph{Output.}
Với mỗi bộ test, in đáp án modulo $10^9+7$.

\paragraph{Sample Input.}
\begin{verbatim}
2
1
2
\end{verbatim}

\paragraph{Sample Output.}
\begin{verbatim}
0
8
\end{verbatim}

\section{1002. I love tree}

\paragraph{Giới hạn.} Thời gian: 1 giây. Bộ nhớ: 512 MB.

Cho một cây có $n$ đỉnh và $q$ thao tác.  Có hai loại thao tác:
\begin{itemize}
  \item \texttt{1 a b}: xét đường đi $\langle a,b\rangle$.  Với đỉnh thứ $x$
  trên đường đi này, cộng thêm $x^2$ vào giá trị của đỉnh đó.  Ví dụ đường đi
  từ $a$ đến $b$ là $(x_1,x_2,x_3,x_4,x_5)$, sau thao tác ta có
  $x_1{+}=1$, $x_2{+}=4$, $x_3{+}=9$, $x_4{+}=16$, $x_5{+}=25$.
  \item \texttt{2 x}: hỏi giá trị của đỉnh $x$.
\end{itemize}

\paragraph{Input.}
Chỉ có một bộ test.  Dòng đầu chứa $n$ ($1\le n\le 10^5$).  Trong $n-1$ dòng
tiếp theo, mỗi dòng chứa hai số $u,v$, biểu diễn một cạnh giữa $u$ và $v$.
Dòng tiếp theo chứa $q$ ($1\le q\le 10^5$).  Mỗi dòng trong $q$ dòng tiếp theo
có dạng \texttt{1 a b} hoặc \texttt{2 x}, với $1\le a,b,x\le n$.

\paragraph{Output.}
Mỗi truy vấn loại 2 in một số nguyên là đáp án.

\paragraph{Sample Input.}
\begin{verbatim}
3
1 2
2 3
5
1 1 2
2 1
1 2 1
2 2
2 3
\end{verbatim}

\paragraph{Sample Output.}
\begin{verbatim}
1
5
0
\end{verbatim}

\section{1003. I love playing games}

\paragraph{Giới hạn.} Thời gian: 1 giây. Bộ nhớ: 512 MB.

Alice và Bob chơi trên một đồ thị vô hướng $G$ có $n$ đỉnh đánh số từ $1$ đến
$n$ và $m$ cạnh.  Ban đầu Alice ở đỉnh $x$, Bob ở đỉnh $y$, và cả hai đều muốn
đến đỉnh $z$.

Alice đi trước, sau đó hai người luân phiên đi.  Trong một lượt, người chơi có
thể đứng yên hoặc đi sang một đỉnh kề.  Hai người không được ở cùng một đỉnh
cùng thời điểm, ngoại trừ vị trí xuất phát và vị trí kết thúc.  Cuối một lượt,
nếu Alice đến $z$ mà Bob chưa đến thì Alice thắng; nếu Bob đến $z$ mà Alice
chưa đến thì Bob thắng; nếu cả hai cùng đến hoặc cùng chưa đến thì kết quả vẫn
có thể là hòa.  Cả hai chơi tối ưu: nếu thắng được thì thắng, nếu không thì cố
gắng hòa.

\paragraph{Input.}
Dòng đầu chứa $T$ ($T\le 10000$).  Mỗi bộ test có $m+2$ dòng.  Dòng đầu chứa
$n,m$ ($n\le 10^3$).  Dòng thứ hai chứa $x,y,z$.  Trong $m$ dòng tiếp theo,
mỗi dòng chứa một cạnh vô hướng $u,v$ ($1\le u,v\le n$).  Tổng $n$ trên mọi bộ
test không vượt quá $10^5$, tổng $m$ không vượt quá $2\cdot 10^5$.  Đồ thị có
thể có khuyên và đa cạnh.

\paragraph{Output.}
Với mỗi bộ test, in $1$ nếu Alice thắng, $2$ nếu hòa, và $3$ nếu Bob thắng.

\paragraph{Sample Input.}
\begin{verbatim}
2
10 13
8 1 10
1 2
1 3
2 4
2 5
3 5
8 9
9 4
9 5
4 6
5 6
6 10
5 7
7 10
8 9
2 3 8
1 2
2 6
2 7
6 4
7 5
4 3
5 3
6 8
7 8
\end{verbatim}

\paragraph{Sample Output.}
\begin{verbatim}
1
1
\end{verbatim}

\section{1004. I love counting}

\paragraph{Giới hạn.} Thời gian: 1 giây. Bộ nhớ: 256 MB.

Cho một dãy độ dài $n$.  Mỗi vị trí có trọng số $c_i$ với $c_i\le n$.  Có $Q$
truy vấn; mỗi truy vấn cho đoạn $(l,r)$ và hai tham số $a,b$.  Hãy đếm số loại
trọng số khác nhau $c$ xuất hiện trong đoạn đó sao cho
\[
  c\oplus a \le b,
\]
trong đó $\oplus$ là phép xor bit.

\paragraph{Input.}
Chỉ có một bộ test.  Dòng đầu chứa $n$ ($n\le 100000$).  Dòng thứ hai chứa
$n$ số $c_i$ ($1\le c_i\le n$).  Dòng thứ ba chứa $Q$ ($Q\le 100000$).  Mỗi
dòng trong $Q$ dòng tiếp theo chứa $l,r,a,b$ với
$1\le l\le r\le n$, $a\le n$, $b\le n$.

\paragraph{Output.}
Với mỗi truy vấn, in số trọng số thỏa điều kiện.

\paragraph{Sample Input.}
\begin{verbatim}
5
1 2 2 4 5
4
1 3 1 3
2 4 4 2
1 5 2 3
4 5 3 6
\end{verbatim}

\paragraph{Sample Output.}
\begin{verbatim}
2
1
2
1
\end{verbatim}

\section{1005. I love string}

\paragraph{Giới hạn.} Thời gian: 1 giây. Bộ nhớ: 512 MB.

Ông X có một chuỗi thao tác, viết thành một xâu.  Ở mỗi thao tác, ký tự tiếp
theo của chuỗi thao tác được chèn vào trước hoặc sau xâu hiện tại.  Ví dụ chuỗi
thao tác là \texttt{aabac}; nếu sau bốn thao tác đầu thu được \texttt{baaa},
thì sau thao tác cuối xâu có thể trở thành \texttt{baaac} hoặc \texttt{cbaaa}.
Thao tác đầu chỉ có một cách, các thao tác còn lại mỗi thao tác có hai cách.

Mỗi cách thao tác có một điểm số: xâu cuối cùng càng nhỏ theo thứ tự từ điển
thì điểm càng cao.  Với chuỗi thao tác cho trước, hãy đếm số cách thao tác đạt
điểm cao nhất.  Hai cách thao tác khác nhau nếu tồn tại một thao tác không phải
thao tác đầu mà một cách chèn ký tự vào trước xâu hiện tại, cách kia chèn vào
sau.

\paragraph{Input.}
Dòng đầu chứa $T$ ($T\le 10$).  Với mỗi bộ test, dòng đầu chứa $n$
($1\le n\le 100000$), dòng tiếp theo chứa một xâu chữ cái thường độ dài $n$.

\paragraph{Output.}
Với mỗi bộ test, in số cách, modulo $1000000007$.

\paragraph{Sample Input.}
\begin{verbatim}
1
5
abcde
\end{verbatim}

\paragraph{Sample Output.}
\begin{verbatim}
1
\end{verbatim}

\section{1006. I love sequences}

\paragraph{Giới hạn.} Thời gian: 1 giây. Bộ nhớ: 512 MB.

Cho ba dãy
$a=[a_1,a_2,\ldots,a_n]$, $b=[b_1,b_2,\ldots,b_n]$ và
$c=[c_1,c_2,\ldots,c_n]$, mỗi dãy gồm $n$ số nguyên không âm.  Cần tính
\[
  \sum_{p=1}^{n}\sum_{k=1}^{n} d_{p,k}\,c_p^k.
\]

Dãy $d=[d_{p,1},d_{p,2},\ldots,d_{p,n}]$ được tính như sau:
\[
  d_{p,k}=\sum_{k=i\oplus j} a_i b_j,
\]
với $1\le p,k\le n$, phạm vi của $i,j$ là
$1\le i\le n/p$, $1\le j\le n/p$.

Ký hiệu $\oplus$ ở đây là phép toán trên biểu diễn tam phân: mỗi chữ số của
$K=I\oplus J$ bằng ước chung lớn nhất của hai chữ số tương ứng trong $I$ và
$J$.  Cụ thể, nếu
$(K)_{10}=(k_{m-1}\ldots k_0)_3$,
$(I)_{10}=(i_{m-1}\ldots i_0)_3$ và
$(J)_{10}=(j_{m-1}\ldots j_0)_3$, thì
$k_t=\gcd(i_t,j_t)$ với mọi $t$.  Quy ước
$\gcd(x,0)=\gcd(0,x)=x$.  Ví dụ nếu
$x=(5)_{10}=(012)_3$ và $y=(10)_{10}=(101)_3$, thì
$x\oplus y=(111)_3=(13)_{10}$.

\paragraph{Input.}
Chỉ có một bộ test.  Dòng đầu chứa $n$ ($n\le 2\cdot 10^5$).  Ba dòng tiếp
theo lần lượt chứa $a_i$, $b_i$, $c_i$ với $1\le a_i,b_i,c_i\le 10^9$.

\paragraph{Output.}
In đáp án modulo $10^9+7$.

\paragraph{Sample Input.}
\begin{verbatim}
4
2 3 4 5
1 2 3 4
9 8 7 6
\end{verbatim}

\paragraph{Sample Output.}
\begin{verbatim}
1643486
\end{verbatim}

\section{1007. I love data structure}

\paragraph{Giới hạn.} Thời gian: 3 giây. Bộ nhớ: 512 MB.

Cần duy trì một dãy số; mỗi vị trí có hai tham số $a$ và $b$.  Có $m$ thao tác,
gồm bốn loại:
\begin{enumerate}
  \item Cho đoạn $(l,r)$, số $x$ và cờ $tag\in\{0,1\}$.  Với mọi vị trí trong
  đoạn, cộng $x$ vào tham số $a$ hoặc $b$ tùy theo cờ.
  \item Cho đoạn $(l,r)$.  Với mỗi vị trí trong đoạn, biến đổi
  $a$ thành $3a+2b$ và $b$ thành $3a-2b$.  Ví dụ từ $a=1,b=2$ thành
  $a=7,b=-1$.
  \item Cho đoạn $(l,r)$.  Với mỗi vị trí trong đoạn, hoán đổi hai tham số
  $a,b$.
  \item Cho đoạn $(l,r)$.  Hỏi
  \[
    \sum_{i=l}^{r} a_i b_i.
  \]
\end{enumerate}

\paragraph{Input.}
Chỉ có một bộ test.  Dòng đầu chứa $n$ ($n\le 200000$).  Trong $n$ dòng tiếp
theo, dòng thứ $i$ chứa $a_i,b_i$ với $1\le a_i,b_i\le 1000000000$.  Dòng tiếp
theo chứa $q$ ($q\le 200000$).  Các thao tác có dạng:
\begin{itemize}
  \item \texttt{1 tag l r x}, với $0\le tag\le 1$, $1\le x\le 1000000000$;
  \item \texttt{2 l r};
  \item \texttt{3 l r};
  \item \texttt{4 l r}.
\end{itemize}

\paragraph{Output.}
Với mỗi truy vấn loại 4, in kết quả modulo $1000000007$.

\paragraph{Sample Input.}
\begin{verbatim}
5
1 4
3 2
4 1
3 6
7 3
6
1 0 2 4 2
2 2 4
4 1 4
3 3 5
1 1 4 5 3
4 1 4
\end{verbatim}

\paragraph{Sample Output.}
\begin{verbatim}
614
623
\end{verbatim}

\section{1008. I love exam}

\paragraph{Giới hạn.} Thời gian: 1 giây. Bộ nhớ: 128 MB.

Sinh viên Z có $n$ môn thi và còn $t$ ngày để ôn tập.  Ban đầu điểm của mọi môn
là $0$.  Bạn học đưa cho Z $m$ bộ tài liệu ôn tập.  Bộ thứ $i$ thuộc môn
$s_i$, cần học $y_i$ ngày và sau đó tăng $x_i$ điểm cho môn đó; điểm mỗi môn
tối đa là $100$, nên khi đã đạt $100$ thì không tăng thêm.  Mỗi bộ tài liệu chỉ
được dùng một lần.

Z không thể học hết tài liệu.  Anh ta cần chọn một số bộ để học sao cho trượt
không quá $p$ môn trong học kỳ này; một môn có điểm dưới $60$ được xem là
trượt.  Hãy tìm tổng điểm lớn nhất có thể đạt được.  Nếu không thể thỏa điều
kiện, in $-1$.

\paragraph{Input.}
Dòng đầu chứa $T$ ($T\le 10$).  Với mỗi bộ test: dòng đầu chứa $n$ ($n\le 50$);
dòng thứ hai chứa $n$ tên môn, mỗi tên dài không quá $15$; dòng thứ ba chứa
$m$ ($m\le 15000$).  Trong $m$ dòng tiếp theo, mỗi dòng chứa tên môn $s$, điểm
tăng $x$ và số ngày $y$ với $1\le x,y\le 10$.  Dữ liệu bảo đảm môn đó có trong
học kỳ này.  Dòng cuối chứa $t,p$ với $1\le t\le 500$, $0\le p\le 3$.

\paragraph{Output.}
Với mỗi bộ test, in tổng điểm lớn nhất thỏa điều kiện, hoặc $-1$ nếu không thể.

\paragraph{Sample Input.}
\begin{verbatim}
1
3
mathematics physics signals
20
physics 10 1
physics 10 1
physics 10 1
physics 10 1
physics 10 1
physics 10 1
physics 10 1
mathematics 10 1
mathematics 10 1
mathematics 10 1
mathematics 10 1
mathematics 10 1
mathematics 10 1
mathematics 10 1
signals 10 1
signals 10 1
signals 10 1
signals 10 1
signals 10 1
signals 10 2
19 1
\end{verbatim}

\paragraph{Sample Output.}
\begin{verbatim}
190
\end{verbatim}

\section{1009. I love triples}

\paragraph{Giới hạn.} Thời gian: 1 giây. Bộ nhớ: 256 MB.

Cho một mảng $a$ gồm $n$ số nguyên.  Hãy đếm số bộ ba $(i,j,k)$ thỏa
$i<j<k$ và $a_i a_j a_k$ là số chính phương.  Số chính phương là số có thể viết
thành tích của hai số nguyên bằng nhau; ví dụ $1,4,9,16,25,36$ là số chính
phương, còn $2,3,6$ thì không.

\paragraph{Input.}
Dòng đầu chứa $T$ ($T\le 6$).  Mỗi bộ test gồm hai dòng: dòng đầu chứa $n$
($1\le n\le 10^5$), dòng thứ hai chứa $n$ số $a_i$ ($1\le a_i\le 10^5$).

\paragraph{Output.}
Với mỗi bộ test, in số bộ ba thỏa điều kiện.

\paragraph{Sample Input.}
\begin{verbatim}
1
6
1 2 4 8 16 32
\end{verbatim}

\paragraph{Sample Output.}
\begin{verbatim}
10
\end{verbatim}

\section{1010. I love permutation}

\paragraph{Giới hạn.} Thời gian: 1 giây. Bộ nhớ: 64 MB.

Cho số nguyên dương $a$ và số nguyên tố lẻ $P$ thỏa $a<P$.  Tạo dãy
\[
  b_x=ax\bmod P,\qquad 1\le x\le P-1,
\]
có độ dài $P-1$.  Hỏi trong dãy này có bao nhiêu cặp nghịch thế.  Vì đáp án có
thể rất lớn, chỉ cần in đáp án modulo $2$.  Một cặp nghịch thế là cặp
$(i,j)$ thỏa $i<j$ và $b_i>b_j$.

\paragraph{Input.}
Dòng đầu chứa $T$ ($T\le 10^5$).  Mỗi bộ test chứa hai số $a,P$ với
$1\le a<P\le 10^{18}$.

\paragraph{Output.}
Với mỗi bộ test, in đáp án của bộ test đó.

\paragraph{Sample Input.}
\begin{verbatim}
4
2 7
3 7
4 7
5 7
\end{verbatim}

\paragraph{Sample Output.}
\begin{verbatim}
0
1
0
1
\end{verbatim}

\section{1011. I love max and multiply}

\paragraph{Giới hạn.} Thời gian: 1 giây. Bộ nhớ: 64 MB.

Cho hai dãy $A_i$ và $B_i$ độ dài $n$, với $0\le i\le n-1$.  Định nghĩa mảng
$C$ độ dài $n$:
\[
  C_k=\max\{A_iB_j\mid (i\mathbin{\&}j)\ge k\},
\]
trong đó $\&$ là phép AND bit.  Hãy tính
\[
  \sum_{i=0}^{n-1} C_i
\]
modulo $998244353$.

\paragraph{Input.}
Dòng đầu chứa $T$.  Mỗi bộ test gồm ba dòng.  Dòng đầu chứa $n$
($1\le n\le 2^{18}$).  Dòng thứ hai chứa $n$ số $a_0,a_1,\ldots,a_{n-1}$ với
$|a_i|\le 10^9$.  Dòng thứ ba chứa $n$ số $b_0,b_1,\ldots,b_{n-1}$ với
$|b_i|\le 10^9$.  Tổng $n$ trên mọi bộ test không vượt quá $2^{19}$.

\paragraph{Output.}
Với mỗi bộ test, in
$ans=\sum_{i=0}^{n-1}C_i \bmod 998244353$.

\paragraph{Sample Input.}
\begin{verbatim}
1
4
9 1 4 1
5 4 1 1
\end{verbatim}

\paragraph{Sample Output.}
\begin{verbatim}
54
\end{verbatim}

\section{1012. I love 114514}

\paragraph{Giới hạn.} Thời gian: 1 giây. Bộ nhớ: 64 MB.

Ông I rất thích \texttt{114514}.  Cho một xâu; nếu \texttt{114514} là xâu con
liên tiếp của xâu đó thì Xiaoi cũng thích xâu này.  Hãy kiểm tra điều đó: nếu
Xiaoi thích xâu, in \texttt{AAAAAA}, ngược lại in \texttt{Abuchulaile}.

\paragraph{Input.}
Dòng đầu chứa $T$ ($T\le 10$).  Mỗi bộ test chứa một xâu $s$ với
$1\le |s|\le 10^5$.

\paragraph{Output.}
Với mỗi bộ test, in đáp án trên một dòng.

\paragraph{Sample Input.}
\begin{verbatim}
2
114514
1919810
\end{verbatim}

\paragraph{Sample Output.}
\begin{verbatim}
AAAAAA
Abuchulaile
\end{verbatim}

\section*{Ghi chú về bản dịch}

Văn bản trích xuất từ PDF bị mất nhiều ký hiệu toán do mã hóa font.  Các ký
hiệu và ràng buộc trong bản dịch này được đối chiếu lại từ trang PDF đã render.

\end{document}
